\documentclass[10pt,a4paper]{ctexart}
\usepackage[margin=2cm]{geometry}
\usepackage{graphicx}
\usepackage{float}      % 支持 [H] 浮动体选项
\usepackage{hyperref}
\usepackage{amsmath}    % 提供方程环境
\usepackage{physics}    % 简化向量、微分算子语法
\newcommand{\myspace}[1]{\par\vspace{#1\baselineskip}} %自定义空一行
\usepackage{fancyhdr}
\pagestyle{fancy}
\fancyhf{}  % 清除所有页眉页脚
\fancyfoot[C]{\thepage}  % 在页脚中央设置页码
\renewcommand{\headrulewidth}{0pt}  % 去掉页眉横线
\renewcommand{\footrulewidth}{0pt}  % 去掉页脚横线

\usepackage{hyperref}
\hypersetup{
    colorlinks=true,
    linkcolor=blue,
    filecolor=magenta,      
    urlcolor=cyan,
    pdftitle={Overleaf Example},
    pdfpagemode=FullScreen,
    }

\usepackage{titlesec}

% 重新定义section格式
\titleformat{\section}
  {\normalfont\Large\bfseries}
  {\chinese{section}、}
  {0pt}
  {}

% 给出常用字号的自定义指令
\newcommand{\chuhao}{\fontsize{42pt}{48pt}\selectfont}
\newcommand{\xiaochu}{\fontsize{36pt}{42pt}\selectfont}
\newcommand{\xiaoer}{\fontsize{18pt}{21.6pt}\selectfont}
\newcommand{\sanhao}{\fontsize{16pt}{20pt}\selectfont}
\newcommand{\xiaosan}{\fontsize{15pt}{18pt}\selectfont}
\newcommand{\sihao}{\fontsize{14pt}{18pt}\selectfont}
\newcommand{\xiaosi}{\fontsize{12pt}{16pt}\selectfont}
\newcommand{\wuhao}{\fontsize{10.5pt}{12.6pt}\selectfont}

% 定义中文数字转换
\titleformat{\section}
  {\normalfont\Large\bfseries}
  {\zhnum{section}、}  % 使用\zhnum而不是\chinese
  {0pt}
  {}

% 标题设置
\title{\xiaochu\textbf{物理实验预习报告}}
\author{}
\date{}

\begin{document}

\vspace*{36pt} % 顶部空白

% 居中填写区域
\begin{center}   
    \xiaochu\textbf{物理实验预习报告}
    % 预留空白区域
    \vspace{13pt}
    \vspace{13pt}
    \vspace{13pt}
    \vspace{13pt}
    
    \sanhao\textbf{实验名称：\underline{\makebox[10cm]{ 光的衍射}}} \\[10pt]
    \sanhao\textbf{指导教师：\underline{\makebox[10cm]{ 姜柏旭 }}} \\[30pt]
    
    % 预留空白区域
    \myspace{6}
    
   \sihao\textbf{班级：\underline{\makebox[6cm]{}}} \\[10pt]
    \sihao\textbf{姓名：\underline{\makebox[6cm]{}}} \\[10pt]
    \sihao\textbf{学号：\underline{\makebox[6cm]{}}} \\[30pt]
    
   \sihao\textbf{实验日期: 
   \underline{\makebox[0.8cm]{2025}}  年 
   \underline{\makebox[0.4cm]{11}}月
   \underline{\makebox[0.4cm]{26}}日  
   星期\underline{\makebox[0.6cm]{三}}上午}
    
    \vspace{18pt}
    \sihao 浙江大学物理实验教学中心
    
\end{center}

\newpage

    \section{实验综述}
    
    \xiaosi （自述实验现象、实验原理和实验方法，不超过300字，5分）

    \subsection{实验原理}
\subsubsection{惠更斯-菲涅尔原理}
因为同一波前上的各点发出的都是相干子波。而各子波在空间某点的相干叠加，就决定了该点波的强度。
\[
E(P) = \int_s {E_{(p)}} = \int_s F\frac{K(\theta)}{r}\cos(\omega t - \frac{2 \pi r}{\lambda})\d{S}
\]
其中$F$是比例系数，$K(\theta)$为倾斜因子。

\subsubsection{单缝衍射光强分布特征}
由菲涅尔原理和矢量振幅法，可知\[
I = I_0 (\frac{\sin \mu}{\mu})^2
\]
想要得到光强度的极值位置，则对$\mu$求偏导
\[
\frac{\partial I}{\partial \mu} = 0
\]
即明纹位置：($a$为狭缝的宽度)
\[
\sin\theta = 0, \pm 1.43\frac{\lambda}{a}, \pm 2.46\frac{\lambda}{a}, \cdots
\]
暗纹位置：
\[
\sin\theta = \pm \frac{\lambda}{a}, \pm \frac{2\lambda}{a}, \cdots
\]
$\pm 1$级明纹的最大光强则为$I_1 \approx 4.7\% I_0$



\subsubsection{利用光栅衍射法测量激光波长}
光栅方程（主极大条件）：
\[
d\sin \theta = \pm k\lambda(k = 0, \pm 1, \pm 2, \cdots)
\]
角度较小时，我们有$\sin\theta \approx \tan \theta \approx \frac{x}{D}$，即
\[
x_k = k\frac{D}{d}\lambda
\]
最后由公式\[\frac{x_1 - x_{-1}}{2} \approx \frac{D}{d}\lambda\]
可以求出激光波长
\subsection{实验方法}
\begin{enumerate}
    \item \textbf{光路调节与校准：}
    调节光学导轨上的各个光学元件，确保它们的位置共轴且等高，以形成一条清晰、水平的实验光路。

    \item \textbf{观察不同孔隙的衍射图样：}
    将具有不同形状孔隙的衍射板置于光路中，定性观察并记录各自产生的夫琅禾费衍射图样的形状特征。

    \item \textbf{利用一维光栅测定激光波长：}
    使用一维光栅进行衍射实验。精确测量其衍射图样中正一级与负一级主极大亮纹之间的距离，然后依据光栅方程 $\frac{x_1 - x_{-1}}{2} \approx \frac{D}{d}\lambda$ 计算出所用激光的波长 $\lambda$。

    \item \textbf{研究单缝衍射的光强分布并测量缝宽：}
    观察由单缝产生的夫琅禾费衍射条纹。利用一维光强测量仪测出其光强分布。绘制出相对光强 $I/I_0$ 随位置坐标 $x$ 变化的函数图像，并根据衍射图样中暗纹的位置，计算出该单缝的精确宽度 $a$。

\end{enumerate}
\myspace{2}
    
    \section{实验重点}
    
    \xiaosi（简述本实验的学习重点，不超过100字，3分）

\begin{enumerate}
    \item \textbf{分析衍射图案：}首先要搞清楚，当光线通过不同形状的物体时，产生的衍射花样的明暗分布有什么不同和规律。
    \item \textbf{准确测量：}准确测量一些光的衍射相关物理量
\end{enumerate}
    \myspace{2}
    
    \section{实验难点}

    \xiaosi（简述本实验的实现难点，不超过100字，2分）
\begin{enumerate}
    \item 实验开始前需要调整光学导轨，确保各光学元件等高
    \item 准确判断衍射图样形状
    \item 转动轮毂时只能单向转动
    \item 测量的时候不能移动光学仪器的位置
\end{enumerate}
\myspace{2}


\newpage

\sihao\textbf{注意事项：}
\xiaosi
\begin{enumerate}
    \item 用PDF格式上传“预习报告”，文件名：学生姓名+学号+实验名称+周次。
    \item “预习报告”必须递交在“学在浙大”的本课程的对应实验项目的“作业”模块内。
    \item “预习报告”还须拷贝到“实验报告”中（便以教师批改）。
    \item 大学物理实验”课程选做实验使用本“预习报告”；必做实验无须写预习报告，前往“学在浙大”完成预习测试即可。
\end{enumerate}

\vspace{1cm}
\xiaosi\centering \textbf{浙江大学物理实验教学中心制}

\end{document}
